\documentclass[12pt,a4paper]{article}

\usepackage{iftex}

\ifPDFTeX
  \usepackage{cmap}
  \usepackage[T2A]{fontenc}
  \usepackage[utf8]{inputenc}
  \usepackage{lmodern}
  \usepackage[english,russian]{babel}
\else
  \usepackage{fontspec}
  \defaultfontfeatures{Ligatures=TeX}
  \setmainfont{Times New Roman}
  \setsansfont{Arial}
  \setmonofont{Consolas}
  \usepackage{polyglossia}
  \setmainlanguage{russian}
  \setotherlanguage{english}
\fi

\usepackage{amsmath,amssymb,amsthm,mathtools}
\usepackage[a4paper,margin=2.2cm,top=2.4cm,bottom=2.2cm]{geometry}
\usepackage{enumitem}
\usepackage{fancyhdr}
\usepackage[hidelinks]{hyperref}

\providecommand{\SheetNo}{}
\newcommand{\sheetauthor}{}

\newcommand{\sheettitle}[1]{
  \begin{center}
    {\Large\bfseries Серия \SheetNo. #1}\\[0.4em]
    \rule{\linewidth}{0.4pt}
  \end{center}
  \vspace{0.5em}
  \markboth{#1}{#1}
}

\setlength{\headheight}{15pt}
\pagestyle{fancy}
\fancyhf{}
\fancyhead[L]{\small Серия \SheetNo}
\fancyhead[R]{\small\leftmark}
\fancyfoot[C]{\thepage}
\renewcommand{\headrulewidth}{0.4pt}

\newlist{problems}{enumerate}{2}
\setlist[problems]{label=\arabic*., ref=\arabic*,
                   leftmargin=*, itemsep=0.8em, topsep=0.8em}

\newcommand{\cyralphnum}[1]{
  \ifcase#1\or а\or б\or в\or г\or д\or е\or ж\or з\or и\or к\or л\or м\or
  н\or п\or р\or с\or т\or у\or ф\or х\or ц\or ч\or ш\or щ\or э\or ю\or я\fi}
\newcommand{\cyralph}[1]{\cyralphnum{\value{#1}}}
\AddEnumerateCounter{\cyralph}{\cyralphnum}{щ}

\newlist{parts}{enumerate}{2}
\setlist[parts]{label=\cyralph*), ref=\cyralph*,
                leftmargin=1.5em, itemsep=0.25em, topsep=0.25em}

\newcommand{\N}{\mathbb{N}}
\newcommand{\Z}{\mathbb{Z}}
\newcommand{\Q}{\mathbb{Q}}
\newcommand{\R}{\mathbb{R}}
\providecommand{\C}{}
\renewcommand{\C}{\mathbb{C}}
\providecommand{\NOD}{}
\renewcommand{\NOD}{\operatorname{\text{НОД}}}
\providecommand{\NOK}{}
\renewcommand{\NOK}{\operatorname{\text{НОК}}}

\def\SheetNo{1}

\begin{document}
\sheettitle{Интро}
\begin{problems}

\item Квадрат $40\times40$ клеток разбит на $400$ L-тетрамино. Докажите, что
найдется прямая, идущая по линии сетки, которая хотя бы $6$ из этих фигур
разрезает на две фигурки из двух клеток.

\item На шахматной доске отметили $25$ клеток. Докажите, что можно расставить
$8$ не бьющих друг друга ладей так, чтобы не менее $4$ ладей стояли на
отмеченных клетках.

\item В стране дураков имеют хождение монеты достоинством $1$, $2$, $4$, $\ldots$,
$2^n$ тугриков. Докажите, что, имея по одной монете каждого достоинства, можно
выплатить любую сумму, не превышающую $2^{n+1}-1$ тугриков.

\item Из клетчатого квадрата $2^n\times 2^n$ вырезали:
\begin{parts}
\item угловую;
\item произвольную клеточку.
\end{parts}
Докажите, что оставшуюся фигуру можно разрезать на уголки из трех клеточек.

\item Плоскость разбита на части несколькими:
\begin{parts}
\item прямыми;
\item окружностями.
\end{parts}
Докажите, что эти части можно раскрасить в 2 цвета правильным образом (т.е. так,
чтобы никакие две области одного цвета не имели общей границы).

\item На плоскости нарисован многоугольник, у которого наружу растут волосы.
Внутри многоугольника проведено несколько диагоналей, на каждой из которых в
одну из сторон растут волосы. Диагонали делят многоугольник на несколько частей.
Докажите, что есть часть, внутри которой нет волос.

\item На столе лежит \(100\) камней. Два игрока по очереди берут из кучки от \(1\) до \(9\) камней. Пропускать ход нельзя. Выигрывает тот, кто забирает последний камень.
Кто из игроков выиграет при правильной игре?

\item На столе лежат \(2021\) красных и \(2022\) зелёных камня. Аня и Петя делают ходы по очереди. 
Аня ходит первой. При каждом ходе игрок выбирает цвет и удаляет $n$ камней этого цвета, где число $n$ должно быть делителем текущего числа камней другого цвета. 
Кто возьмёт последний камень, тот выиграет. Кто из игроков может обеспечить себе победу независимо от ходов соперника?

\item На трибунах хоккейной арены несколько рядов по $168$ мест в каждом ряду.
На финальный матч в качестве зрителей пригласили $2016$ учеников нескольких
спортивных школ, не более $40$ от каждой школы. Учеников любой школы требуется
разместить на одном ряду. Какое наименьшее количество рядов должно быть на арене,
чтобы это в любом случае удалось сделать?

\item В клетчатом квадрате $11\times11$ отметили все $144$ вершины клеток. Затем
отмеченные точки раскрасили в пять цветов. При каком наибольшем $d$ могло оказаться,
что расстояние между любыми двумя одноцветными отмеченными точками не меньше $d$?

\end{problems}
\end{document}
